% !TEX encoding = UTF-8 Unicode
\documentclass[hyperref={bookmarks=false}]{beamer}
\usetheme{CambridgeUS}
\setbeamercolor{bibliography entry author}{fg=black}
\setbeamercolor{item}{fg=darkred}
\setbeamercolor{caption name}{fg=darkred}
\beamertemplatenavigationsymbolsempty

\usepackage[serbianc]{babel}
\usepackage{type1ec}
\usepackage{cmap}
\usepackage{tikz}

\title{\textbf{\color{darkred}МАТ\color{black}Ф\color{darkred}ЕТЕРИЈА}}
\titlegraphic{\includegraphics[height=3cm,width=3.5cm]{square-matf-light.png}}
\author{Лазар Васовић \and Душан Пантелић}
\institute[]{Математички факултет, Универзитет у Београду\\\href{https://gitlab.com/matfpveb/projekti/2019-2020/36-matfeterija}{gitlab.com/matfpveb/projekti/2019-2020/36-matfeterija}}
\date[Математички факултет]{31. август 2020.}

\begin{document}

\frame{\titlepage}

\begin{frame}{Садржај}
\tableofcontents[subsectionstyle=hide]
\end{frame}

\section{Опште}
\subsection{Функционалности}
\begin{frame}{Функционалности}
\begin{itemize}
\item Студентска огласна табла Математичког факултета

\item Регистрација, потврда преко мејла, пријава

\item Преглед, додавање, лајковање, брисање објава

\item Преглед, додавање, лајковање, брисање коментара

\item Филтрирање по већем броју критеријума

\item Подешавање свог налога, преглед и праћење других

\item Информација о тренутно активним корисницима

\item Прототипна соба за госте -- анонимне поруке

\end{itemize}
\end{frame}

\section{Клијент}
\subsection{Клијент}
\begin{frame}{Клијент}
\begin{itemize}
\item meAn -- Angular, TypeScript

\item Дизајн -- одговоран, у духу актуелног визуелног идентитета факултета, неколико анимација и графичких занимљивости

\item Све са часа -- све врсте директива, сервиса и везивања података, референце, рутирање, филтери, реактивни формулари, HTTP...

\item Пресретачи -- допуна захтева креденцијалима, обрада грешака

\item Робусност -- кликтање дугмади, рад при паду базе или сервера, ауторизациони сервис (canActivate), рад са претплатама...

\item Ажурност -- превазилажење тога што је HTTP без стања

\item Соба за госте -- vanilla Js, fetch API, манипулација DOM-ом
\end{itemize}
\end{frame}

\section{Сервер}
\subsection{Сервер}
\begin{frame}{Сервер}
\begin{itemize}
\item mEaN -- Node.js, Express.js

\item Све са часа --  путање и датотеке, рутирање, функције средњег слоја, креативни EJS за излиставање интерфејса...

\item REST API -- JSON одговори и по једна датотека и HTML

\item Аватари -- сервер их чува, док је у бази само линк за GET

\item Безбедност и аутентификација -- JWT колачић, RSA енкрипција, потврда налога преко званичног факултетског мејла

\end{itemize}
\end{frame}

\section{База}
\subsection{База}
\begin{frame}{База}
\begin{itemize}
\item Mean -- MongoDB, Mongoose ODM

\item Реплицирана нерелациона база --  дистрибуиране трансакције, редундантност података, већи одзив и доступност

\item Све са часа -- схеме, обавезне и подразумеване вредности, низови, страни кључеви и њихово попуњавање, разни упити...

\item Индекси -- обични, сложени, јединствени, временски (TTL)

\item Окидачи над схемама -- решавање каскадних зависности

\item Безбедност -- лозинке заштићене хеширањем и селектором

\item Соба за госте -- Monk, једноставно али неструктурирано
\end{itemize}
\end{frame}

\section{}
\subsection{Захвалница}
\begin{frame}
\centering \LARGE
\textbf{ХВАЛА НА ПАЖЊИ!}

\textbf{Питања?}
\end{frame}

\end{document}